% =====================================================================================
% Vision Mamba for Diabetic Retinopathy Grading -- project report
%
% Build (from the report/ folder):
%     pdflatex main && bibtex main && pdflatex main && pdflatex main
% or simply:
%     latexmk -pdf main.tex
%
% Figures are read from report/figures/, produced by `python -m eval.plots`.
% If a figure does not exist yet, a labelled placeholder box is drawn instead, so
% the report always compiles.
%
% Every place that needs numbers from the real training run is marked with
%     % TODO: fill in after training
% =====================================================================================
\documentclass[11pt,a4paper]{article}

\usepackage[utf8]{inputenc}
\usepackage[T1]{fontenc}
\usepackage[margin=2.5cm]{geometry}
\usepackage{amsmath,amssymb}
\usepackage{graphicx}
\usepackage{booktabs}
\usepackage{algorithm}
\usepackage{algpseudocode}
\usepackage[numbers]{natbib}
\usepackage{xcolor}
\usepackage[hidelinks]{hyperref}

\graphicspath{{figures/}}

% \maybefigure{file}{width}: include the figure if it exists, otherwise draw a
% framed placeholder naming the command that creates it.
\newcommand{\maybefigure}[2]{%
  \IfFileExists{figures/#1}{%
    \includegraphics[width=#2]{#1}%
  }{%
    \fbox{\parbox[c][4cm][c]{#2}{\centering\small\color{gray}
      Figure \texttt{#1} not generated yet.\\
      Run \texttt{python -m eval.plots} after training.}}%
  }%
}

% Placeholder marker for numbers that only exist after training
\newcommand{\todo}{\textcolor{gray}{--}}

\title{Vision Mamba for Diabetic Retinopathy Grading:\\
       a From-Scratch Implementation and Cross-Dataset Study}
\author{Marc Pedrin\\IISc Spectrum Lab internship task}
\date{\today}

\begin{document}
\maketitle

% -------------------------------------------------------------------------------------
\begin{abstract}
We implement Vision Mamba (Vim)~\cite{zhu2024vim}, a vision backbone built from
bidirectional selective state space models, from scratch in plain PyTorch, and apply
it to five-class diabetic retinopathy (DR) grading. Every component, including the
selective scan, is written explicitly; no Mamba kernels, pretrained weights or
ready-made data loaders are used. The model is trained on EyePACS and evaluated both
on a held-out, patient-disjoint EyePACS test split and, without any adaptation, on
the external Messidor-2 dataset. We use a scaled-down configuration (128-pixel
inputs, 12 blocks) suited to a single GPU, and document every deviation from the
paper. % TODO: fill in after training -- one sentence with the headline QWK on both
% datasets, e.g. "The model reaches a quadratic weighted kappa of X on EyePACS and Y
% on Messidor-2, a drop we attribute to domain shift."
\end{abstract}

% -------------------------------------------------------------------------------------
\section{Introduction}

Diabetic retinopathy is damage to the retina caused by diabetes, and a leading cause
of preventable blindness in working-age adults. It progresses silently, so
diabetic patients are screened regularly with fundus photographs, colour images of
the back of the eye. Screening at population scale needs many trained graders,
which motivates automated grading.

\paragraph{Grading scale.} We use the five-level international clinical scale
\cite{wilkinson2003scale}, as in the EyePACS and Messidor-2 labels:
\begin{center}
\begin{tabular}{cll}
\toprule
Grade & Name & Typical findings \\
\midrule
0 & No DR & no abnormalities \\
1 & Mild non-proliferative & microaneurysms only \\
2 & Moderate non-proliferative & more than microaneurysms, less than grade 3 \\
3 & Severe non-proliferative & extensive haemorrhages, venous beading, IRMA \\
4 & Proliferative & new vessel growth and/or vitreous haemorrhage \\
\bottomrule
\end{tabular}
\end{center}
The grades are \emph{ordinal}: confusing grade~0 with grade~4 is a much worse error
than confusing grade~1 with grade~2. The lesions that separate the low grades are
tiny (a microaneurysm spans a few pixels at our resolution), and the classes are
heavily imbalanced, with about 73\% of EyePACS images at grade~0. Both facts shape
our preprocessing, our loss and our choice of metric.

\paragraph{Why Vision Mamba.} Vision Transformers~\cite{dosovitskiy2021vit} relate
every image patch to every other through self-attention, at a cost quadratic in the
number of patches. Vim replaces attention with selective state space models
(Mamba~\cite{gu2023mamba}), whose cost is linear in the sequence length, and makes
them bidirectional so that every patch can receive information from both sides.

\paragraph{Contributions of this project.}
(i) A line-by-line, from-scratch implementation of the Vim block (Algorithm~1 of the
paper), tested against hand-computed cases; (ii) a leakage-free, patient-level
EyePACS pipeline with fundus-specific preprocessing; (iii) from-scratch metrics,
including quadratic weighted kappa, verified against scikit-learn; (iv) an in-domain
and cross-domain evaluation.

% -------------------------------------------------------------------------------------
\section{The Vision Mamba architecture}

\subsection{State space models}
A state space model maps an input signal $x(t)$ to an output $y(t)$ through a hidden
state $h(t)\in\mathbb{R}^N$:
\[
  h'(t) = \mathbf{A}h(t) + \mathbf{B}x(t), \qquad y(t) = \mathbf{C}h(t).
\]
This continuous form is theory only and is never coded. To process a sequence of
tokens it is discretised with a step size $\Delta$, using the zero-order hold:
\begin{equation}
  \overline{\mathbf{A}} = \exp(\Delta\mathbf{A}), \qquad
  \overline{\mathbf{B}} = (\Delta\mathbf{A})^{-1}\bigl(\exp(\Delta\mathbf{A})-\mathbf{I}\bigr)
                          \cdot \Delta\mathbf{B},
  \label{eq:zoh}
\end{equation}
which gives the recurrence
\begin{equation}
  h_t = \overline{\mathbf{A}}\,h_{t-1} + \overline{\mathbf{B}}\,x_t, \qquad
  y_t = \mathbf{C}\,h_t .
  \label{eq:recurrence}
\end{equation}
When $\overline{\mathbf{A}}$, $\overline{\mathbf{B}}$ and $\mathbf{C}$ are the same at
every step (as in S4~\cite{gu2022s4}), the recurrence can be unrolled into a single
global convolution with kernel
\begin{equation}
  \overline{\mathbf{K}} = \bigl(\mathbf{C}\overline{\mathbf{B}},\;
     \mathbf{C}\overline{\mathbf{A}}\,\overline{\mathbf{B}},\; \dots,\;
     \mathbf{C}\overline{\mathbf{A}}^{M-1}\overline{\mathbf{B}}\bigr),
  \qquad y = x * \overline{\mathbf{K}} .
  \label{eq:conv}
\end{equation}
Equation numbers (\ref{eq:zoh})--(\ref{eq:conv}) coincide with the paper's
Eqs.~1--3.

\paragraph{Selectivity.} Mamba makes $\mathbf{B}$, $\mathbf{C}$ and $\Delta$
\emph{functions of the current token}, computed by small linear layers. Each token
can then decide how much to write into the state, what to read out of it, and how
much of the past to forget (a large $\Delta$ drives $\overline{\mathbf{A}}$ towards
zero). Because these quantities change at every step, no fixed kernel
$\overline{\mathbf{K}}$ exists and Eq.~(\ref{eq:conv}) no longer applies: the
recurrence (\ref{eq:recurrence}) must be run, which we do with an explicit loop.
$\mathbf{A}$ is diagonal and kept strictly negative, so
$\overline{\mathbf{A}}=\exp(\Delta\mathbf{A})\in(0,1)$ and the state decays rather
than explodes.

\subsection{From images to token sequences}
An image of size $H\times W$ is cut into $J$ patches of $P\times P$ pixels, each
linearly projected to $D$ dimensions. A learnable class token $t_{\mathrm{cls}}$ is
added and learnable position embeddings $\mathbf{E}_{\mathrm{pos}}$ are summed in:
\begin{equation}
  \mathbf{T}_0 = [\,t_{\mathrm{cls}};\; t_p^1\mathbf{W};\; t_p^2\mathbf{W};\;\dots;\;
                 t_p^J\mathbf{W}\,] + \mathbf{E}_{\mathrm{pos}} .
  \label{eq:tokens}
\end{equation}
($t_{\mathrm{cls}}$ is printed at the front, but the paper's default model places it
in the middle of the sequence; see Section~\ref{sec:ablations}.) The sequence passes
through $L$ Vim blocks with residual connections,
$\mathbf{T}_l = \mathrm{Vim}(\mathbf{T}_{l-1}) + \mathbf{T}_{l-1}$, and the class
token's final output is normalised and classified:
$f = \mathrm{Norm}(\mathbf{T}_L^0)$, $\hat p = \mathrm{MLP}(f)$.

\subsection{The bidirectional Vim block}
Algorithm~\ref{alg:vim} reproduces the paper's Algorithm~1, with the same line
numbering (comment lines included). Our \texttt{vim/vim\_block.py} and
\texttt{vim/selective\_scan.py} follow it line by line, and each line of code is
annotated with the algorithm line it implements.

\begin{algorithm}[t]
\caption{Vim block process (after Algorithm~1 of \citet{zhu2024vim}).}
\label{alg:vim}
\small
\begin{algorithmic}[1]
\Require token sequence $\mathbf{T}_{l-1}$ : $(\mathtt{B},\mathtt{M},\mathtt{D})$
\Ensure token sequence $\mathbf{T}_{l}$ : $(\mathtt{B},\mathtt{M},\mathtt{D})$
\State \textit{/* normalise the input sequence $\mathbf{T}'_{l-1}$ */}
\State $\mathbf{T}'_{l-1} : (\mathtt{B},\mathtt{M},\mathtt{D}) \gets \mathbf{Norm}(\mathbf{T}_{l-1})$
\State $\mathbf{x} : (\mathtt{B},\mathtt{M},\mathtt{E}) \gets \mathbf{Linear}^{\mathbf{x}}(\mathbf{T}'_{l-1})$
\State $\mathbf{z} : (\mathtt{B},\mathtt{M},\mathtt{E}) \gets \mathbf{Linear}^{\mathbf{z}}(\mathbf{T}'_{l-1})$
\State \textit{/* process with different direction */}
\For{$o$ in \{forward, backward\}}
  \State $\mathbf{x}'_o : (\mathtt{B},\mathtt{M},\mathtt{E}) \gets \mathbf{SiLU}(\mathbf{Conv1d}_o(\mathbf{x}))$
  \State $\mathbf{B}_o : (\mathtt{B},\mathtt{M},\mathtt{N}) \gets \mathbf{Linear}^{\mathbf{B}}_o(\mathbf{x}'_o)$
  \State $\mathbf{C}_o : (\mathtt{B},\mathtt{M},\mathtt{N}) \gets \mathbf{Linear}^{\mathbf{C}}_o(\mathbf{x}'_o)$
  \State \textit{/* softplus ensures positive $\Delta_o$ */}
  \State $\Delta_o : (\mathtt{B},\mathtt{M},\mathtt{E}) \gets \log\bigl(1+\exp(\mathbf{Linear}^{\Delta}_o(\mathbf{x}'_o) + \mathbf{Parameter}^{\Delta}_o)\bigr)$
  \State \textit{/* shape of $\mathbf{Parameter}^{\mathbf{A}}_o$ is $(\mathtt{E},\mathtt{N})$ */}
  \State $\overline{\mathbf{A}}_o : (\mathtt{B},\mathtt{M},\mathtt{E},\mathtt{N}) \gets \Delta_o \otimes \mathbf{Parameter}^{\mathbf{A}}_o$
  \State $\overline{\mathbf{B}}_o : (\mathtt{B},\mathtt{M},\mathtt{E},\mathtt{N}) \gets \Delta_o \otimes \mathbf{B}_o$
  \State \textit{/* initialise $\mathbf{h}_o$ and $\mathbf{y}_o$ with 0 */}
  \State $\mathbf{h}_o : (\mathtt{B},\mathtt{E},\mathtt{N}) \gets \mathrm{zeros}(\mathtt{B},\mathtt{E},\mathtt{N})$
  \State $\mathbf{y}_o : (\mathtt{B},\mathtt{M},\mathtt{E}) \gets \mathrm{zeros}(\mathtt{B},\mathtt{M},\mathtt{E})$
  \State \textit{/* SSM recurrent */}
  \For{$i$ in $\{0, \dots, M-1\}$}
    \State $\mathbf{h}_o = \overline{\mathbf{A}}_o[:,i,:,:] \odot \mathbf{h}_o + \overline{\mathbf{B}}_o[:,i,:,:] \odot \mathbf{x}'_o[:,i,:,\mathrm{None}]$
    \State $\mathbf{y}_o[:,i,:] = \mathbf{h}_o \otimes \mathbf{C}_o[:,i,:]$
  \EndFor
\EndFor
\State \textit{/* get gated $\mathbf{y}_o$ */}
\State $\mathbf{y}'_{\mathrm{forward}} : (\mathtt{B},\mathtt{M},\mathtt{E}) \gets \mathbf{y}_{\mathrm{forward}} \odot \mathbf{SiLU}(\mathbf{z})$
\State $\mathbf{y}'_{\mathrm{backward}} : (\mathtt{B},\mathtt{M},\mathtt{E}) \gets \mathbf{y}_{\mathrm{backward}} \odot \mathbf{SiLU}(\mathbf{z})$
\State \textit{/* residual connection */}
\State $\mathbf{T}_l : (\mathtt{B},\mathtt{M},\mathtt{D}) \gets \mathbf{Linear}^{\mathbf{T}}(\mathbf{y}'_{\mathrm{forward}} + \mathbf{y}'_{\mathrm{backward}}) + \mathbf{T}_{l-1}$
\end{algorithmic}
\end{algorithm}

Line~13 is printed without the exponential of Eq.~(\ref{eq:zoh}); we apply it, as the
official code does, since otherwise $\overline{\mathbf{A}}$ would be negative.
Line~14 is the simplified discretisation $\overline{\mathbf{B}}=\Delta\mathbf{B}$,
a first-order approximation of Eq.~(\ref{eq:zoh}) that is also used by the official
Mamba code. The backward direction is implemented by flipping the sequence, running
the same code with the backward weights, and flipping the result back.

\subsection{Design choices validated by the paper's ablations}
\label{sec:ablations}
\paragraph{Bidirectionality (paper Table~4).} A single forward scan only carries
information left to right, so a token never hears from the patches after it.
\begin{center}
\begin{tabular}{lcc}
\toprule
Bidirectional strategy & ImageNet top-1 (\%) & ADE20K mIoU \\
\midrule
None (forward only)                     & 73.2 & 32.3 \\
Bidirectional layer                     & 70.9 & 33.6 \\
Bidirectional SSM                       & 72.8 & 33.2 \\
Bidirectional SSM + Conv1d (default)    & \textbf{73.9} & \textbf{35.9} \\
\bottomrule
\end{tabular}
\end{center}

\paragraph{Class token position (paper Table~5).} Because the scans are directional
and their memory decays with distance, a class token in the middle of the sequence
is reached by both directions with the shortest average distance to every patch.
\begin{center}
\begin{tabular}{lc}
\toprule
Classification design & ImageNet top-1 (\%) \\
\midrule
Mean pool               & 73.9 \\
Max pool                & 73.4 \\
Head class token        & 75.2 \\
Double class token      & 74.3 \\
Middle class token (default) & \textbf{76.1} \\
\bottomrule
\end{tabular}
\end{center}
We use the default variant in both cases: bidirectional SSM with Conv1d, and a
middle class token at index 32 of the 65-token sequence.

% -------------------------------------------------------------------------------------
\section{Implementation details and deviations from the paper}

Everything is written from scratch in PyTorch: patch embedding, tokens and position
embeddings, the selective scan (an explicit Python loop over the sequence), the
bidirectional block, the model, the datasets, the split and the metrics. PyTorch
supplies only generic building blocks (\texttt{nn.Linear}, \texttt{nn.Conv1d},
\texttt{nn.LayerNorm}, autograd, AdamW) and torchvision only image transforms.
The paper's hardware-aware parallel scan (its Section~3.5) is deliberately out of
scope.

\begin{center}
\begin{tabular}{lll}
\toprule
Setting & Paper (Vim-Ti) & Ours \\
\midrule
Input size            & $224\times224$ & $128\times128$ \\
Patch size / patches  & 16 / 196 & 16 / 64 \\
Sequence length $M$   & 197 & 65 \\
Embedding dim $D$     & 192 & 192 \\
Expanded dim $E$      & 384 & 384 \\
State dim $N$         & 16 & 16 \\
Depth $L$             & 24 & 12 \\
Conv1d kernel         & 4 & 4 \\
Class token           & middle & middle (index 32) \\
$\Delta$ projection   & low-rank (code) & full $E\times E$ (as printed) \\
Norm                  & RMSNorm (code) & LayerNorm \\
$D\cdot x$ skip term  & yes (code) & no (not in Algorithm~1) \\
Parameters            & $\approx$7M & 6,857,093 \\
\bottomrule
\end{tabular}
\end{center}

\paragraph{Parameter count.} Halving the depth does not halve the parameter count,
because the printed Algorithm~1 projects $\Delta$ with a full $E\times E$ matrix
(147,456 weights per direction per block), where the official code uses a rank-12
factorisation. Our block has 557,952 parameters, roughly twice an official block,
and the two $\Delta$ projections account for 53\% of it.

\paragraph{Initialisation.} $\mathbf{A}=-\exp(\mathbf{A}_{\log})$ with
$\mathbf{A}_{\log}$ initialised to $\log(1,\dots,16)$ (S4D-real~\cite{gu2022s4d}),
giving a spread of decay rates; the $\Delta$ bias is set so that the initial step
sizes are log-uniform in $[0.001, 0.1]$, as in Mamba; the class token and position
embeddings use a truncated normal with standard deviation 0.02.
We observed that, with these small initial step sizes and no $D\cdot x$ term, a
freshly initialised block adds very little to its input (standard deviation about
$4\times10^{-4}$ for a unit-variance input), so the initial 12-block stack is close
to the identity. This is stable, and AdamW's per-parameter step normalisation lets
training grow the branches.

\paragraph{Efficiency of the reference loop.} Indexing the precomputed
$\overline{\mathbf{A}}$ inside the loop makes the backward pass allocate a full-size
gradient at every step. Splitting the tensors once with \texttt{torch.unbind}
leaves the mathematics unchanged and made the backward pass about $11\times$ faster
on CPU.

\paragraph{Verification.} 63 unit and integration tests cover, among others: a
two-step selective scan computed by hand; the limits of instant forgetting and of
no forgetting (a running sum); separate parameters per direction and a control
experiment showing the block is direction-aware; gradient flow through 12 blocks;
patient-disjoint splits; and every metric against scikit-learn.

% -------------------------------------------------------------------------------------
\section{Datasets and preprocessing}

\paragraph{EyePACS} (Kaggle Diabetic Retinopathy Detection, 2015~\cite{eyepacs2015}):
35,126 labelled images from 17,563 patients (left and right eye each), from many
camera models. We split it at the \emph{patient} level into 80\% train, 10\%
validation and 10\% test, stratified by each patient's worse-eye grade, with a fixed
seed. Keeping both eyes of a patient in the same split prevents data leakage: the
two eyes are highly correlated, and a model could otherwise score well partly by
recognising the patient.

\paragraph{Messidor-2}~\cite{decenciere2014messidor,abramoff2013messidor2}:
1,748 images from 874 examinations, acquired in French hospitals with a different
camera, with the adjudicated grades of \citet{krause2018grader}. Ungradable images
are excluded. Messidor-2 is used \emph{only} for evaluation; it never influences
training, checkpoint selection or any hyperparameter.

\paragraph{Preprocessing.} (1) Crop the black border to the bounding box of the
retina (pixels brighter than 10/255); (2) pad to a square with black so the round
retina is not distorted; (3) resize to $128\times128$; (4) normalise with the
ImageNet channel statistics. Training images are additionally flipped horizontally
and vertically, rotated by up to $\pm15^\circ$, and jittered by up to 10\% in
brightness, contrast and saturation (no hue change, since lesion colour is
diagnostic). We deliberately avoid the paper's random crops, random erasing and
mixup: they can remove or blend away the few tiny lesions that define the lower
grades, silently making the label wrong.

% -------------------------------------------------------------------------------------
\section{Training methodology}

\begin{center}
\begin{tabular}{p{3.6cm}p{4.2cm}p{6.4cm}}
\toprule
Hyperparameter & Paper (ImageNet) & Ours (EyePACS) \\
\midrule
Optimiser          & AdamW ($\beta_1=0.9$) & AdamW ($\beta=(0.9, 0.999)$) \\
Learning rate      & $1\times10^{-3}$ & $1\times10^{-3}$ \\
Weight decay       & 0.05 & 0.05 (not on norms, biases, $\mathbf{E}_{\mathrm{pos}}$, $t_{\mathrm{cls}}$, $\mathbf{A}_{\log}$) \\
Schedule           & cosine & 3 warmup epochs, cosine to $10^{-5}$ \\
Epochs             & 300 & 30 \\
Batch size         & 1024 & 32 \\
Loss               & CE + label smoothing & CE with inverse-frequency class weights \\
Augmentation       & crop, flip, mixup, erasing & flips, rotation $\pm15^\circ$, mild jitter \\
EMA                & yes & no \\
Gradient clipping  & -- & max norm 1.0 \\
Precision          & -- & mixed (float16 autocast) \\
Checkpoint selection & -- & best validation QWK \\
Seed               & -- & 42 (all generators, deterministic cuDNN) \\
Hardware           & 8 A800 GPUs & 1 Kaggle T4 \\
\bottomrule
\end{tabular}
\end{center}

The short schedule and small batch reflect a limited GPU budget; the goal of the
project is a faithful, understandable implementation rather than a state-of-the-art
score. Class weights $w_c = N/(K\,n_c)$ counter the imbalance, so that each grade
contributes about equally to the loss.

\paragraph{Metrics.} We report accuracy, macro-averaged precision, recall and F1, and
quadratic weighted kappa (QWK)~\cite{cohen1968weighted}, the official EyePACS
metric:
\[
  \kappa = 1 - \frac{\sum_{i,j} W_{ij} O_{ij}}{\sum_{i,j} W_{ij} E_{ij}},\qquad
  W_{ij} = \frac{(i-j)^2}{(K-1)^2},
\]
where $O$ is the confusion matrix and $E$ the matrix expected by chance from the
marginal grade frequencies. The quadratic weights punish a prediction four grades
away sixteen times as hard as a prediction one grade away, matching the ordinal
nature of DR grades. Accuracy alone is misleading here: always predicting grade~0
already scores about 73\%, but has $\kappa=0$.

% -------------------------------------------------------------------------------------
\section{Results}

% TODO: fill in after training -- copy the numbers printed by
% `python -m eval.evaluate` (also stored in outputs/default/results.json).
\begin{table}[h]
\centering
\caption{Results on the held-out EyePACS test split (in-domain) and on Messidor-2
(external, never seen in training).}
\label{tab:results}
\begin{tabular}{lcc}
\toprule
Metric & EyePACS test & Messidor-2 \\
\midrule
Images           & \todo & \todo \\
Accuracy         & \todo & \todo \\
QWK              & \todo & \todo \\
Macro precision  & \todo & \todo \\
Macro recall     & \todo & \todo \\
Macro F1         & \todo & \todo \\
\bottomrule
\end{tabular}
\end{table}

% TODO: fill in after training -- per-grade recall from the evaluate.py table
\begin{table}[h]
\centering
\caption{Per-grade recall (share of each true grade that was predicted correctly).}
\label{tab:recall}
\begin{tabular}{lcc}
\toprule
Grade & EyePACS test & Messidor-2 \\
\midrule
0 none      & \todo & \todo \\
1 mild      & \todo & \todo \\
2 moderate  & \todo & \todo \\
3 severe    & \todo & \todo \\
4 proliferative & \todo & \todo \\
\bottomrule
\end{tabular}
\end{table}

\begin{figure}[h]
\centering
\maybefigure{training_curves.png}{\textwidth}
\caption{Convergence: training and validation loss and accuracy, and validation QWK
per epoch; the marked point is the checkpoint used for evaluation.}
\label{fig:curves}
\end{figure}
% TODO: fill in after training -- two or three sentences describing Fig. 1:
% does the loss still fall at epoch 30, is there a train/validation gap
% (over-fitting), and at which epoch did validation QWK peak?

\begin{figure}[h]
\centering
\maybefigure{confusion_eyepacs_test.png}{0.48\textwidth}\hfill
\maybefigure{confusion_messidor2.png}{0.48\textwidth}
\caption{Row-normalised confusion matrices on the EyePACS test split (left) and
Messidor-2 (right). Cell colour is the share of the true grade; the count is in
brackets.}
\label{fig:confusion}
\end{figure}
% TODO: fill in after training -- which grades are confused most? Are errors mostly
% to neighbouring grades (good for QWK) or far away? Is grade 1 recognised at all?

\begin{figure}[h]
\centering
\maybefigure{examples_eyepacs_test.png}{\textwidth}\\[1ex]
\maybefigure{examples_messidor2.png}{\textwidth}
\caption{Example predictions on EyePACS (top) and Messidor-2 (bottom): true and
predicted grade, half chosen among errors.}
\label{fig:examples}
\end{figure}

% -------------------------------------------------------------------------------------
\section{Analysis of cross-dataset generalisation}

A drop in performance from the EyePACS test split to Messidor-2 is expected, even
though the task is the same, because of \emph{domain shift}: Messidor-2 was acquired
with a different camera (colour balance, sharpness, field of view and border
shape differ), in a different population (French hospital patients rather than a
mostly US screening network), under different acquisition protocols, and graded to
a different, adjudicated reference standard. Whatever the model learnt that is
specific to EyePACS photographs rather than to the lesions themselves does not
transfer. Our preprocessing (cropping to the retina, padding, normalisation, colour
jitter) is designed to reduce the most obvious of these differences, but cannot
remove them.

% TODO: fill in after training -- quantify the drop (QWK and per-grade recall from
% Tables 1-2), say which grades suffer most, and relate it to the confusion matrices
% (e.g. does the model under-call grades on the external set?). Compare with the
% grade distributions of the two datasets.

% -------------------------------------------------------------------------------------
\section{Limitations}

\begin{itemize}
  \item \textbf{Scale.} 128-pixel inputs discard detail: microaneurysms become a
        pixel or two, which particularly hurts grade~1. Twelve blocks and 30 epochs
        are far below the paper's 24 blocks and 300 epochs.
  \item \textbf{Speed.} The selective scan is a sequential Python loop, much slower
        than the paper's fused parallel kernel; this limited the training budget.
  \item \textbf{No pretraining.} The model is trained from random initialisation on
        about 28,000 images, whereas Vim-Ti is trained on 1.28 million ImageNet
        images.
  \item \textbf{Label noise.} EyePACS grades come from single graders and are known
        to be noisy, which caps the achievable agreement.
  \item \textbf{Single run.} Results are from one seed; we did not measure the
        variance across seeds.
        % TODO: fill in after training -- remove or adjust if several seeds are run
  \item \textbf{Faithfulness choices.} Following the printed Algorithm~1 (full
        $\Delta$ projection, no $D\cdot x$ term) makes our blocks larger than the
        official ones and our model a slightly different function from the official
        Vim-Ti.
\end{itemize}

% -------------------------------------------------------------------------------------
\section{Conclusion}

We implemented Vision Mamba from first principles, keeping a direct, documented
correspondence between the paper's equations and Algorithm~1 and the code, and
applied it to diabetic retinopathy grading with a leakage-free evaluation on an
in-domain test split and an external dataset.
% TODO: fill in after training -- summarise the headline QWK on both datasets and
% the size of the domain-shift drop, in one or two sentences.
The main lessons are that the selective, bidirectional scan can be implemented
transparently in a few dozen lines, that faithful choices such as the printed full
$\Delta$ projection have measurable consequences for model size, and that
cross-dataset evaluation is essential for judging a medical imaging model.

\bibliographystyle{plainnat}
\bibliography{references}

\end{document}
